% GENERATED FILE. Edit canonical lessons, then run npm run generate:books.
% canonical-source-hash: fnv1a64:40309f4f31a1d5a3
% canonical-lessons: TE-C343-ksatagatrulu, TE-C343-durghatana, TE-C343-tokkisalata, TE-C343-trutilo, TE-C343-bayatapadu

\chapter{A Narrow Escape}
\label{ch:te-a-narrow-escape}
\hlchaptermodality{343}

\begin{chapteropening}
\textbf{By the end of this chapter:} \emph{I can say the injured, a tragic accident, a stampede, narrowly and to get out.}

Complete the last lesson of chapter 343: I can say the injured, a tragic accident, a stampede, narrowly and to get out.
\end{chapteropening}

\section[\texorpdfstring{\te{క్షతగాత్రులు} (kṣatagātrulu)}{క్షతగాత్రులు (kṣatagātrulu)}]{\te{క్షతగాత్రులు} (kṣatagātrulu) — the injured}

\label{lesson:TE-C343-ksatagatrulu}



\begin{quote}
Before the new one: say the Telugu for a speech, then the Telugu for to give a speech.
\end{quote}

\subsection*{You'll want to know: \te{క్షతగాత్రులు}}
\textbf{\te{క్షతగాత్రులు}} — \emph{kṣatagātrulu} — \textquotedblleft{}the injured\textquotedblright{}.

\begin{quote}
\te{క్షతగాత్రులను} \te{ఆసుపత్రికి} \te{తీసుకెళ్ళారు}.
\end{quote}

\begin{quote}
\emph{kṣatagātrulanu āsupatriki tīsukeḷḷāru.}
\end{quote}

\begin{quote}
The injured were taken to hospital.
\end{quote}

A news word, nearly only in the plural: \te{క్షత}, wounded, and \te{గాత్ర}, the body. At the hospital the \te{క్షతగాత్రులు} are given \te{ఔషధం}.

\begin{cousinweb}
\textbf{\te{క్షతగాత్రులు}} (\emph{kṣatagātrulu}) is borrowed from Sanskrit: \emph{kṣata}, \textquotedblleft{}wounded\textquotedblright{} (CDIAL 3645), and \emph{gātra}, \textquotedblleft{}a limb, the body\textquotedblright{} (CDIAL 4124): those wounded in body.
\end{cousinweb}

\begin{grammarlens}[title={What you've built}]
One more word for telling what happened.
\end{grammarlens}

\subsection*{Your turn}
\begin{itemize}
\raggedright
  \item \emph{Say it:} \emph{kṣatagātrulu}
  \item \emph{Say it:} \emph{kṣatagātrulu}, once more
  \item \emph{Say it:} \emph{kṣatagātrulanu āsupatriki tīsukeḷḷāru.} — The injured were taken to hospital.
  \item \emph{Recall:} say the Telugu for a speech, then the Telugu for to give a speech, then say \emph{kṣatagātrulu} again
\end{itemize}

\begin{tcolorbox}[breakable,colback=teal!4,colframe=teal!35!black,title={Before you move on}]
What does \te{క్షతగాత్రులు} mean? (\textquotedblleft{}The injured\textquotedblright{}.) Say it once more.
\end{tcolorbox}
\section[\texorpdfstring{\te{దుర్ఘటన} (durghaṭana)}{దుర్ఘటన (durghaṭana)}]{\te{దుర్ఘటన} (durghaṭana) — a tragic accident; a disaster}

\label{lesson:TE-C343-durghatana}



\begin{quote}
Before the new one: say the Telugu for to give a speech, then the Telugu for the injured.
\end{quote}

\subsection*{You'll want to know: \te{దుర్ఘటన}}
\textbf{\te{దుర్ఘటన}} — \emph{durghaṭana} — \textquotedblleft{}a tragic accident; a disaster\textquotedblright{}.

\begin{quote}
\te{ఆ} \te{దుర్ఘటన} \te{నిన్న} \te{సాయంత్రం} \te{జరిగింది}.
\end{quote}

\begin{quote}
\emph{ā durghaṭana ninna sāyantraṁ jarigindi.}
\end{quote}

\begin{quote}
That tragic accident happened yesterday evening.
\end{quote}

\te{దుర్}-, ill, with \te{ఘటన}, a happening: compare \te{సంఘటన}.

\begin{cousinweb}
\textbf{\te{దుర్ఘటన}} (\emph{durghaṭana}) is built from Sanskrit \emph{dur-}, \textquotedblleft{}ill, bad\textquotedblright{}, and \emph{ghaṭana}, \textquotedblleft{}a joining, a coming together\textquotedblright{} (Monier-Williams).
\end{cousinweb}

\begin{grammarlens}[title={What you've built}]
One more word for telling what happened.
\end{grammarlens}

\subsection*{Your turn}
\begin{itemize}
\raggedright
  \item \emph{Say it:} \emph{durghaṭana}
  \item \emph{Say it:} \emph{durghaṭana}, once more
  \item \emph{Say it:} \emph{ā durghaṭana ninna sāyantraṁ jarigindi.} — That tragic accident happened yesterday evening.
  \item \emph{Recall:} say the Telugu for to give a speech, then the Telugu for the injured, then say \emph{durghaṭana} again
\end{itemize}

\begin{tcolorbox}[breakable,colback=teal!4,colframe=teal!35!black,title={Before you move on}]
What does \te{దుర్ఘటన} mean? (\textquotedblleft{}A tragic accident; a disaster\textquotedblright{}.) Say it once more.
\end{tcolorbox}
\section[\texorpdfstring{\te{తొక్కిసలాట} (tokkisalāṭa)}{తొక్కిసలాట (tokkisalāṭa)}]{\te{తొక్కిసలాట} (tokkisalāṭa) — a stampede}

\label{lesson:TE-C343-tokkisalata}



\begin{quote}
Before the new one: say the Telugu for the injured, then the Telugu for a tragic accident.
\end{quote}

\subsection*{You'll want to know: \te{తొక్కిసలాట}}
\textbf{\te{తొక్కిసలాట}} — \emph{tokkisalāṭa} — \textquotedblleft{}a stampede\textquotedblright{}.

\begin{quote}
\te{జాతరలో} \te{తొక్కిసలాట} \te{జరిగింది}.
\end{quote}

\begin{quote}
\emph{jātaralō tokkisalāṭa jarigindi.}
\end{quote}

\begin{quote}
There was a stampede at the fair.
\end{quote}

From \te{తొక్కు}, to tread on: people treading on one another.

\begin{cousinweb}
\textbf{\te{తొక్కిసలాట}} (\emph{tokkisalāṭa}) is native Dravidian, built on \emph{tokku}, \emph{trokku}, \textquotedblleft{}to tread, trample\textquotedblright{} (DEDR 3522).
\end{cousinweb}

\begin{grammarlens}[title={What you've built}]
One more word for telling what happened.
\end{grammarlens}

\subsection*{Your turn}
\begin{itemize}
\raggedright
  \item \emph{Say it:} \emph{tokkisalāṭa}
  \item \emph{Say it:} \emph{tokkisalāṭa}, once more
  \item \emph{Say it:} \emph{jātaralō tokkisalāṭa jarigindi.} — There was a stampede at the fair.
  \item \emph{Recall:} say the Telugu for the injured, then the Telugu for a tragic accident, then say \emph{tokkisalāṭa} again
\end{itemize}

\begin{tcolorbox}[breakable,colback=teal!4,colframe=teal!35!black,title={Before you move on}]
What does \te{తొక్కిసలాట} mean? (\textquotedblleft{}A stampede\textquotedblright{}.) Say it once more.
\end{tcolorbox}
\section[\texorpdfstring{\te{త్రుటిలో} (truṭilō)}{త్రుటిలో (truṭilō)}]{\te{త్రుటిలో} (truṭilō) — narrowly; by a whisker}

\label{lesson:TE-C343-trutilo}



\begin{quote}
Before the new one: say the Telugu for a tragic accident, then the Telugu for a stampede.
\end{quote}

\subsection*{You'll want to know: \te{త్రుటిలో}}
\textbf{\te{త్రుటిలో}} — \emph{truṭilō} — \textquotedblleft{}narrowly; by a whisker\textquotedblright{}.

\begin{quote}
\te{అతను} \te{త్రుటిలో} \te{ప్రమాదం} \te{నుంచి} \te{తప్పించుకున్నాడు}.
\end{quote}

\begin{quote}
\emph{atanu truṭilō pramādaṁ nuñci tappiñcukunnāḍu.}
\end{quote}

\begin{quote}
He narrowly escaped the accident.
\end{quote}

\te{త్రుటి}, a tiny instant, with -\te{లో}, in: in a split second. Headlines often spell it \te{తృటిలో}.

\begin{cousinweb}
\textbf{\te{త్రుటిలో}} (\emph{truṭilō}) is built on Sanskrit \emph{truṭi}, \textquotedblleft{}a small space of time\textquotedblright{} (CDIAL 6064), with the Telugu locative \emph{-lō}, \textquotedblleft{}in\textquotedblright{}.
\end{cousinweb}

\begin{grammarlens}[title={What you've built}]
One more word for telling what happened.
\end{grammarlens}

\subsection*{Your turn}
\begin{itemize}
\raggedright
  \item \emph{Say it:} \emph{truṭilō}
  \item \emph{Say it:} \emph{truṭilō}, once more
  \item \emph{Say it:} \emph{atanu truṭilō pramādaṁ nuñci tappiñcukunnāḍu.} — He narrowly escaped the accident.
  \item \emph{Recall:} say the Telugu for a tragic accident, then the Telugu for a stampede, then say \emph{truṭilō} again
\end{itemize}

\begin{tcolorbox}[breakable,colback=teal!4,colframe=teal!35!black,title={Before you move on}]
What does \te{త్రుటిలో} mean? (\textquotedblleft{}Narrowly; by a whisker\textquotedblright{}.) Say it once more.
\end{tcolorbox}
\section[\texorpdfstring{\te{బయటపడు} (bayaṭapaḍu)}{బయటపడు (bayaṭapaḍu)}]{\te{బయటపడు} (bayaṭapaḍu) — to get out; to come to light}

\label{lesson:TE-C343-bayatapadu}



\begin{quote}
Before the new one: say the Telugu for a stampede, then the Telugu for narrowly.
\end{quote}

\subsection*{You'll want to know: \te{బయటపడు}}
\textbf{\te{బయటపడు}} — \emph{bayaṭapaḍu} — \textquotedblleft{}to get out; to come to light\textquotedblright{}.

\begin{quote}
\te{మంటల} \te{నుంచి} \te{అందరూ} \te{క్షేమంగా} \te{బయటపడ్డారు}.
\end{quote}

\begin{quote}
\emph{maṇṭala nuñci andarū kṣēmaṅgā bayaṭapaḍḍāru.}
\end{quote}

\begin{quote}
Everyone got out of the fire safely.
\end{quote}

Past (I …): \te{బయటపడ్డాను}, \emph{bayaṭapaḍḍānu}.

\te{బయట}, outside, with \te{పడు}: to come out, or for a secret to come out.

\begin{cousinweb}
\textbf{\te{బయటపడు}} (\emph{bayaṭapaḍu}) is a native compound: \emph{bayaṭa}, \textquotedblleft{}outside\textquotedblright{}, which sits beside \emph{bayalu}, \textquotedblleft{}the outside, open space\textquotedblright{} (DEDR 3940), with the helping verb \emph{paḍu}.
\end{cousinweb}

\begin{grammarlens}[title={What you've built}]
One more word for telling what happened.
\end{grammarlens}

\subsection*{Your turn}
\begin{itemize}
\raggedright
  \item \emph{Say it:} \emph{bayaṭapaḍu}
  \item \emph{Say it:} \emph{bayaṭapaḍu}, once more
  \item \emph{Say it:} \emph{maṇṭala nuñci andarū kṣēmaṅgā bayaṭapaḍḍāru.} — Everyone got out of the fire safely.
  \item \emph{Recall:} say the Telugu for a stampede, then the Telugu for narrowly, then say \emph{bayaṭapaḍu} again
\end{itemize}

\begin{tcolorbox}[breakable,colback=teal!4,colframe=teal!35!black,title={Before you move on}]
What does \te{బయటపడు} mean? (\textquotedblleft{}To get out; to come to light\textquotedblright{}.) Say it once more.
\end{tcolorbox}
